% RQ1 Architecture: Simulation Framework

\subsection{System Overview}

Our simulation framework models DAO governance as a discrete-time multi-agent system. The architecture comprises four layers:

\begin{enumerate}
    \item \textbf{Configuration Layer}: YAML experiment definitions specifying parameters
    \item \textbf{Agent Layer}: Heterogeneous agents with participation behaviors
    \item \textbf{Governance Layer}: Proposal lifecycle and voting mechanics
    \item \textbf{Analysis Layer}: Metric collection and statistical export
\end{enumerate}

\subsection{Agent Implementation}

Agents are autonomous decision-makers implementing the participation model from Section~\ref{sec:theory}.

\subsubsection{Agent Types}

The framework supports heterogeneous populations:

\begin{table}[htbp]
\centering
\caption{Agent archetypes for participation modeling}
\label{tab:agent-types}
\begin{tabular}{lcccc}
\toprule
Type & $\theta$ (propensity) & Token Share & Fatigue Sensitivity \\
\midrule
Whale & 0.8 & 20-30\% & Low (0.5$\times$) \\
Delegate & 0.9 & 5-10\% & Medium (1.0$\times$) \\
Active & 0.5 & 30-40\% & Medium (1.0$\times$) \\
Passive & 0.1 & 20-30\% & High (2.0$\times$) \\
\bottomrule
\end{tabular}
\end{table}

\subsubsection{Participation Decision}

Each simulation step, agents evaluate active proposals:

\begin{lstlisting}[caption=Agent participation logic (pseudocode)]
def decide_participation(agent, proposal):
    if agent.fatigue >= agent.fatigue_threshold:
        return ABSTAIN

    relevance = compute_relevance(agent, proposal)
    effective_propensity = agent.theta * (1 - agent.fatigue)

    if random() < effective_propensity * relevance:
        agent.fatigue += FATIGUE_INCREMENT
        return VOTE
    return ABSTAIN
\end{lstlisting}

\subsection{Governance Mechanics}

\subsubsection{Proposal Lifecycle}

Proposals progress through states:
\begin{equation}
\text{Draft} \rightarrow \text{Active} \rightarrow \text{Voting} \rightarrow \{\text{Passed}, \text{Failed}, \text{NoQuorum}\}
\end{equation}

Quorum is evaluated at voting close. A proposal passes if:
\begin{enumerate}
    \item Turnout $\geq$ quorum threshold
    \item Yes votes $>$ No votes (simple majority)
\end{enumerate}

\subsubsection{Quorum Configuration}

The framework supports:
\begin{itemize}
    \item \textbf{Fixed quorum}: Constant percentage threshold
    \item \textbf{Dynamic quorum}: Adjusts based on rolling participation average
    \item \textbf{Participation targets}: Configurable target rates that influence agent behavior
\end{itemize}

\subsection{Metrics Collection}

For participation analysis, we capture:

\begin{description}
    \item[Turnout] Fraction of eligible tokens that voted
    \item[Quorum Rate] Fraction of proposals achieving quorum
    \item[Retention] Agent participation consistency over time
    \item[Fatigue Distribution] Population fatigue statistics per timestep
    \item[Concentration] Gini coefficient of actual voting participation
\end{description}

\subsection{Simulation Loop}

Each step represents one day:

\begin{algorithmic}
\FOR{step = 1 to max\_steps}
    \STATE Generate proposals based on frequency parameter
    \STATE Update agent fatigue (decay toward recovery)
    \FOR{each active voting proposal}
        \FOR{each agent}
            \STATE Agent decides participation
        \ENDFOR
        \STATE Tally votes, check quorum
    \ENDFOR
    \STATE Record turnout, quorum achievement, fatigue stats
\ENDFOR
\end{algorithmic}

\subsection{Reproducibility}

All simulations use deterministic seeding. Configuration hashes identify unique experiments. Results include manifests with git commits, seeds, and verification hashes.
